% CLAUDE REPORT::WCBRIDGE-0001::2026-11-15T01:50Z
% !TEX program = xelatex
% Compiles with XeLaTeX or pdfLaTeX (ASCII source).
\documentclass[11pt]{article}
\usepackage[margin=1in]{geometry}
\usepackage{amsmath,amssymb,amsthm}
\usepackage{xurl}
\newtheorem{theorem}{Theorem}
\newtheorem{proposition}[theorem]{Proposition}
\newtheorem{corollary}[theorem]{Corollary}
\newtheorem{lemma}[theorem]{Lemma}
\theoremstyle{definition}
\newtheorem{definition}[theorem]{Definition}
\theoremstyle{remark}
\newtheorem*{remark}{Remark}
\newcommand{\SAT}{\mathrm{SAT}}
\newcommand{\UNSAT}{\mathrm{UNSAT}}
\newcommand{\Br}{\mathrm{Br}}
\newcommand{\LB}{\mathrm{LB}}
\newcommand{\poly}{\operatorname{poly}}
\newcommand{\status}[1]{\textbf{[#1]}}

\title{CLAUDE WCBRIDGE-0001\\[2pt]
\large The worst-case derivation bridge $\SAT\in\mathrm P\Rightarrow\overline W_R(N)\le\poly(N)$:\\
conservation of difficulty, monotonicity, the Resolution specialization test,\\
and the one known non-vacuous instance}
\author{CLAUDE (scheduled routine, independent research track)}
\date{2026-11-14 17:50 America/Los\_Angeles}

\begin{document}
\maketitle

\noindent\textbf{Status.} P versus NP remains \textbf{OPEN}. This report is additive: it changes no
task schedules, roles, allocations, ownership regions or root \texttt{main.tex} (the CLAUDE track has
no ownership region there). It addresses the shared worst-case objective
\[
\overline W_R(N)=\max_{F\in\UNSAT,\ |\langle F\rangle|\le N}\ \min_{\pi:F\vdash_R\bot}\operatorname{cost}_R(\pi)
\]
and the conditional bridge $\Br(R):\ \SAT\in\mathrm P\Rightarrow\overline W_R(N)\le\poly(N)$.
No computation was run; every claim below is a proof or a cited theorem.

\section{Setting and standing assumption}
$R$ is a sound and complete refutation system for CNF in the sense of Cook and Reckhow: a
polynomial-time checker accepts $(F,\pi)$ only if $F\in\UNSAT$, and every $F\in\UNSAT$ has some
accepted $\pi$. \textbf{Assumption C (cost/length):} there is $c\ge1$ with
$|\pi|^{1/c}-c\le\operatorname{cost}_R(\pi)\le |\pi|^{c}+c$ for every accepted $\pi$. (This holds for
unit node cost on shared DAGs with polynomially bounded line size; it must be audited for the
project tariff of $R_{\mathrm{mix341}}$ before any statement below is applied to it.)
Write $\LB(R)$ for ``$\overline W_R(N)\ge N^{\omega(1)}$''.

\section{The free-choice clarification is correct}
\begin{proposition}[\status{PROVED/EXACT}]
If an algorithm $A$ is $R$-based, i.e.\ on every $F\in\UNSAT$ its run determines an accepted
$R$-refutation $\pi_A(F)$ with $\operatorname{cost}_R(\pi_A(F))\le \mathrm{time}_A(F)$, then
$\max_{|F|\le N,F\in\UNSAT}\mathrm{time}_A(F)\ge\overline W_R(N)$.
\end{proposition}
\begin{proof}$\mathrm{time}_A(F)\ge\operatorname{cost}(\pi_A(F))\ge\min_\pi\operatorname{cost}(\pi)$; take max over $F$.\end{proof}
Granting the best derivation for free therefore loses nothing for the lower bound. The difficulty is
entirely in the hypothesis ``$A$ is $R$-based'', which an arbitrary SAT algorithm does not satisfy.
That gap \emph{is} $\Br(R)$. (Atserias--M\"uller 2019: finding short Resolution refutations is NP-hard;
this concerns search and is irrelevant to the minimum, consistent with the directive.)

\section{Conservation of difficulty}
\begin{theorem}[\status{PROVED/EXACT}, elementary]\label{thm:cons}
For every system $R$:
\begin{enumerate}
\item[(a)] $\mathrm P\ne\mathrm{NP}\Rightarrow\Br(R)$ (vacuously).
\item[(b)] $\Br(R)$ is false iff $\SAT\in\mathrm P$ and $\overline W_R$ is not polynomially bounded.
Hence no counterexample to $\Br(R)$ exists unless one also exhibits (or proves the existence of) a
polynomial-time SAT algorithm.
\item[(c)] If $\LB(R)$ is a theorem, then $\Br(R)\iff\SAT\notin\mathrm P$.
\item[(d)] If $R$ is p-optimal (p-simulates every Cook--Reckhow system), then $\Br(R)$ is a theorem
and $\LB(R)\iff\mathrm{NP}\ne\mathrm{coNP}$.
\end{enumerate}
\end{theorem}
\begin{proof}(a),(b) are propositional. (c): given $\LB(R)$, $\Br(R)$ reads $\SAT\in\mathrm P\Rightarrow\bot$.
(d): $\SAT\in\mathrm P$ gives $\mathrm{NP}=\mathrm{coNP}$, so a polynomially bounded system $Q$ exists
(Cook--Reckhow 1979); $R$ p-simulates $Q$, so $\overline W_R$ is polynomial by Assumption C. The
equivalence in (d) is Cook--Reckhow applied to the p-optimal $R$.\end{proof}

\begin{corollary}[\status{PROVED/EXACT}; consequence of known lower bounds]
For each of the following $R$, proving $\Br(R)$ is \emph{literally equivalent} to proving
$\mathrm P\ne\mathrm{NP}$, so the decomposition ``$\LB(R)$ + $\Br(R)$'' gives no reduction in
difficulty: dag-like Resolution (Haken 1985; random 3-CNF: Chv\'atal--Szemer\'edi 1988,
Ben-Sasson--Wigderson 2001); Polynomial Calculus and PCR (Razborov 1998; Ben-Sasson--Impagliazzo 1999;
Alekhnovich--Razborov 2001); Cutting Planes (Pudlák 1997, via interpolation); bounded-depth Frege
(Ajtai 1988; Kraj\'\i\v{c}ek--Pudl\'ak--Woods 1995; Pitassi--Beame--Impagliazzo 1993);
regular Res($\oplus$) (Efremenko--Garl\'\i k--Itsykson 2024; Bhattacharya--Chattopadhyay--Dvo\v{r}\'ak 2024);
Res($k$) for small $k$.
\end{corollary}

\begin{lemma}[Downward monotonicity, \status{PROVED/EXACT}]
If $R'$ p-simulates $R$ (and both satisfy C), then $\overline W_{R'}\le\poly(\overline W_R)$, hence
$\Br(R)\Rightarrow\Br(R')$ and $\LB(R')\Rightarrow\LB(R)$.
\end{lemma}
So the bridge becomes \emph{weaker} (easier) exactly as the system becomes stronger, and the lower bound
becomes harder at the same rate. At the weak end (Corollary) the bridge is the whole problem; at the
p-optimal end (Theorem~\ref{thm:cons}(d)) the lower bound is $\mathrm{NP}\ne\mathrm{coNP}$, which is
stronger than the goal. \textbf{Any genuine progress must come from a system in the middle} for which
$\LB(R)$ is currently \textbf{OPEN} and $\Br(R)$ is provable by an explicit construction:
candidates are general dag-like Res($\oplus$) (open: Itsykson, ITCS 2026, ECCC TR25-116), AC$^0[p]$-Frege,
Frege, Extended Frege, and the project's $R_{\mathrm{mix341}}$ if its lower bound is unknown.

\section{The one known non-vacuous bridge}
\begin{theorem}[\status{KNOWN THEOREM}: Cook 1975; see Kraj\'\i\v{c}ek 1995, Ch.~9]
Let $A$ be a polynomial-time algorithm, correct on all $F\in\UNSAT$ (i.e.\ $A(F)=0$), such that the
bounded-arithmetic theory PV proves soundness of its rejections:
$\mathrm{PV}\vdash\forall\varphi\,\forall a\,(A(\varphi)=0\to\neg\mathrm{Sat}(\varphi,a))$.
Then Extended Frege is polynomially bounded, so $\overline W_{\mathrm{EF}}(N)\le\poly(N)$.
\end{theorem}
\noindent\emph{Proof sketch.} The propositional translation of the PV theorem has EF proofs of size
$\poly(n)$ for inputs of length $n$; substitute the concrete $F$ for $\varphi$; EF verifies the concrete
polynomial-time computation $A(F)=0$ in polynomial size; modus ponens yields $\neg F(a)$.
This bridge was already adopted in the corpus (LSPV-1/357 ``formal-arithmetic one-policy sufficient
condition''); here it is placed in the conservation picture. Its contrapositive is the strongest
unconditional-hypothesis statement available today: \emph{an EF lower bound shows that no
polynomial-time SAT algorithm is PV-provably sound} (a ``feasible'' P$\ne$NP). The missing theorem is
\[
\textbf{(RR) reflection removal:}\quad \SAT\in\mathrm P\ \Longrightarrow\ \exists A\ \text{poly-time, correct, with }
\mathrm{PV}\vdash\mathrm{Sound}_A,
\]
which is \textbf{OPEN} and not expected to follow from known methods (cf.\ Kraj\'\i\v{c}ek--Pudl\'ak--Takeuti 1991:
PV-provable $\mathrm{P}=\mathrm{NP}$ collapses the bounded-arithmetic hierarchy; Pich--Santhanam 2021).
Status of $\Br(\mathrm{EF})$ without the PV hypothesis: \textbf{OPEN}.

\section{The Resolution specialization test (red-team tool)}
\begin{proposition}[\status{PROVED/EXACT}]
Let $\mathcal T$ be a purported proof of $\Br(R)$ that uses only properties of $R$ that are also held by
dag-like Resolution (soundness, completeness, Assumption C, closure under restriction, any locality of
rules satisfied by Resolution, etc.). Then $\mathcal T$ yields $\Br(\mathrm{Res})$, and with Haken's
theorem it proves $\mathrm P\ne\mathrm{NP}$ outright.
\end{proposition}
Hence every candidate bridge must be checked for the exact property of $R$ it uses that
Resolution (and Cutting Planes, PC, bounded-depth Frege) lacks. If none is named, the argument is either
a complete P$\ne$NP proof in disguise or contains an error; historically it has always been the latter.
In Cook's theorem the distinguishing property is that EF proves the propositional translations of PV
theorems, which Resolution does not.

\section{Hard UNSAT families: survival across systems}
\begin{itemize}
\item \textbf{Tseitin on expanders}: hard for Resolution and bounded-depth Frege, \emph{easy} for
GF(2) reasoning, Res($\oplus$)-tree with Gaussian steps, CP-with-parity, and the corpus's paid
AND3+XOR $R_{\mathrm{mix341}}$ proof of cost $4h-1$. FALSIFIED as a candidate for any system with linear algebra.
\item \textbf{PHP}: hard for Resolution and bounded-depth Frege; polynomial in Cutting Planes and EF. FALSIFIED for strong systems.
\item \textbf{Random 3-CNF, $m=\Delta n$, $\Delta$ above threshold}: exponential for Resolution and PC;
exponential for Res($k$), $k=O(\sqrt{\log n})$ (Carenini, ECCC TR26-158; secondary, to be verified);
only superlinear for depth-$k$ Frege (Gryaznov--Talebanfard 2024, arXiv 2403.02275);
Cutting Planes lower bounds known only for $\Theta(\log n)$-CNFs (Fleming--Pankratov--Pitassi--Robere 2017;
Hrubeš--Pudl\'ak 2017), constant width \textbf{OPEN}. No known polynomial upper bound in any system
(Feige's hypothesis). This is the only standard family that survives every known method and is therefore
the preferred CANDIDATE. Its hardness for any middle system in the sense of Section~3 is \textbf{OPEN}.
\end{itemize}

\section{Barrier-transfer and PQ ledger}
\begin{itemize}
\item Bounded arithmetic $\leftrightarrow$ propositional proofs (Cook 1975; Buss): \textbf{directly transferable theorem} (Section 4).
\item PQ item 6 (special soundness/extraction): \textbf{analogy}; Cook's theorem is an ``extractor'' that needs a provable
correctness witness. Extracting $R$-proofs from a black-box solver alone is a \textbf{failed transfer}: the clocked
system $Q_A$ (codex 2026-11-11, eq.~(18)) has trivially short proofs, and $R$-simulation of $Q_A$ is exactly the
open step.
\item PQ item 5 (worst/average-case): \textbf{analogy/candidate}; random 3-CNF hardness for strong systems is an
average-case statement and does not bridge to arbitrary algorithms (Theorem~\ref{thm:cons}(a)--(b)).
\end{itemize}

\section{Assignments implied for the five roles}
\begin{enumerate}
\item \textbf{Lower-Bound Lead}: state $\Br$ only for a middle system; record which side (bridge or bound) is open, never both as ``weakest''.
\item \textbf{Reduction Builder}: determine whether $R_{\mathrm{mix341}}$ p-simulates EF (unrestricted extension).
If yes, $\LB(R_{\mathrm{mix341}})$ implies an EF lower bound (major open problem); if it is p-simulated by a system
in the Corollary list, $\Br(R_{\mathrm{mix341}})$ is equivalent to P$\ne$NP.
\item \textbf{Hypergraph \& Certificates}: prioritize random 3-CNF; retire Tseitin/PHP as all-system candidates.
\item \textbf{Resource Auditor}: verify Assumption C for the project tariff.
\item \textbf{Red Team}: apply the Resolution specialization test to every bridge claim; note that $\Br(R)$ cannot be
falsified by a counterexample formula (Theorem~\ref{thm:cons}(b)).
\end{enumerate}

\section*{CORPUS REVIEWED}
\texttt{AGENTS.md}; root \texttt{main.tex} ownership markers and $R_{\mathrm{mix341}}$/$\overline W_R$ passages;
\texttt{research/agents/lspv-1/2026-11-13\_USER\_TOTAL\_PROGRESS\_TWO\_GATE\_REFOCUS.md} (Gates S/T);
\texttt{research/agents/codex/2026-11-11\_1331\_TRANSFER\_RECOVERY\_AND\_RATE\_AMPLIFICATION.tex} ($Q_A$, p-optimality);
LSPV-1/357, /361 (Cook 1975 reflection, ER), /392 (pointwise Haken B1 counterexample);
recent TMR-1..5 commit log (MCSP firewall track, not bearing on this bridge); prior CLAUDE reports.

\section*{DUPLICATION / CONTRADICTION AUDIT}
Duplicated: Cook--Reckhow p-optimality consequence and $Q_A$ (codex 11-11); Cook-PV sufficient condition (LSPV-1/357).
New here: Theorem~\ref{thm:cons}(b)--(c) worst-case equivalence $\Br(R)\iff\mathrm{P}\ne\mathrm{NP}$ for every
lower-bounded $R$ (LSPV-1/392 is the pointwise analogue only); downward monotonicity; the Resolution specialization test;
the survival table. No contradiction found.

\section*{WHAT CHANGED RELATIVE TO ROOT MAIN}
Nothing edited. Root \texttt{main.tex} already keeps B3/Gate T OPEN; this report supplies a reason why Gate T
cannot be weaker than P$\ne$NP for any $R$ with a known lower bound.

\section*{ADOPTED / REJECTED / RETIRED PRIOR COMPONENTS}
Adopted: $Q_A$, Cook-PV bridge, Gate S/T split. Rejected as all-system hard candidates: Tseitin, PHP. Retired: none.

\section*{STALE OR OVERSTATED CLAIMS FOUND}
Describing $\Br(R)$ as the ``weakest sufficient bridge'' is accurate only relative to stronger simulation theorems
for the same $R$; across systems it is not weaker than P$\ne$NP whenever $\LB(R)$ is known.

\section*{References}
Cook (STOC 1975); Cook--Reckhow (JSL 1979); Haken (TCS 1985); Chv\'atal--Szemer\'edi (JACM 1988);
Ben-Sasson--Wigderson (JACM 2001); Razborov (1998); Alekhnovich--Razborov (2001); Pudl\'ak (JSL 1997);
Ajtai (1988); Kraj\'\i\v{c}ek, \emph{Bounded Arithmetic, Propositional Logic and Complexity Theory} (1995);
Kraj\'\i\v{c}ek--Pudl\'ak--Takeuti (APAL 1991); Pich--Santhanam (STOC 2021);
Atserias--M\"uller (FOCS 2019); Fleming--Pankratov--Pitassi--Robere (FOCS 2017); Hrube\v{s}--Pudl\'ak (FOCS 2017);
Gryaznov--Talebanfard, arXiv 2403.02275; Itsykson, ITCS 2026 / ECCC TR25-116; Efremenko--Itsykson, CCC 2025 / ECCC TR25-039;
Byramji--Impagliazzo ECCC TR25-118; Carenini ECCC TR26-158; Efremenko--Garl\'\i k--Itsykson (STOC 2024).

\end{document}
% END_OF_VERIFIED_REPORT::CLAUDE::WCBRIDGE-0001
